\section{Introduzione}
\subsection{Scopo del documento}
Il presente documento ha lo scopo di descrivere l'analisi dei requisiti del progetto \textbf{AIRA} (AI Requests Automation), relativo al capitolato C0 proposto dall'azienda Sanmarco Informatica S.p.A. Il documento definisce in modo strutturato gli attori, i casi d'uso e i requisiti del prodotto, classificandoli in obbligatori, desiderabili e opzionali, sulla base del capitolato, degli incontri con la proponente e delle linee guida da essa fornite al gruppo \textbf{13oolean}.

\subsection{Scopo del prodotto}
Il prodotto mira a realizzare uno strumento multimodale basato sull'intelligenza artificiale che riceve, interpreta ed estrae il significato semantico da canali eterogenei, come e-mail con allegati e messaggi vocali, convertendo informazioni non strutturate in dati strutturati. I dati, organizzati in formato JSON, vengono inoltrati automaticamente alle piattaforme gestionali aziendali (portali B2B, ERP, CRM e sistemi di ticketing) tramite chiamate API REST o notifiche e-mail.

In questo modo si riduce la gestione manuale dei flussi informativi, con una diminuzione del rischio di errori e dei tempi di lavorazione delle pratiche commerciali. La prima versione del prodotto (MVP) è focalizzata sull'elaborazione degli ordini d'acquisto B2B ricevuti via e-mail e tramite messaggi vocali WhatsApp, come descritto nella sezione \ref{sec:descrizione}.

\subsection{Glossario}
Al fine di evitare ambiguità, i termini tecnici o specialistici sono definiti nel documento \textit{Glossario}. Come stabilito dalle \textit{Norme di Progetto}, tali termini sono evidenziati in corsivo e contrassegnati con il pedice G alla loro prima occorrenza in ciascuna sezione.

\subsection{Riferimenti}
\subsubsection{Normativi}
	\begin{itemize}
		\item Capitolato C0 - AIRA, AI Requests Automation (Sanmarco Informatica S.p.A.);
		\item Regolamento del Progetto - Corso di Ingegneria del Software A.A. 2026/2027;
		\item \textit{Norme di Progetto}.
	\end{itemize}

\subsubsection{Informativi}
	\begin{itemize}
		\item Slide del corso di Ingegneria del Software, A.A. 2026/2027 (T05 - Analisi dei Requisiti);
		\item \textit{Glossario};
		\item \textit{Verbale Riunione Esterna \#1} del 2026-10-21, tenuta con Sanmarco Informatica;
		\item Email della proponente (Linda Benozzato) contenente le linee guida per lo sviluppo del progetto, ricevuta in data 2026-11-12.
	\end{itemize}

\subsection{Fonti dei requisiti}
Per ciascun requisito il documento indica la fonte da cui è stato individuato. Le fonti utilizzate sono:
	\begin{itemize}
		\item \textbf{Capitolato}: documento di presentazione del capitolato C0;
		\item \textbf{VE-1.n}: punto \textit{n} del resoconto del \textit{Verbale Riunione Esterna \#1} del 2026-10-21;
		\item \textbf{LG-n.m}: punto \textit{m} della sezione \textit{n} dell'email contenente le linee guida della proponente, dove la sezione 1 riguarda perimetro dell'MVP e gestione delle intenzioni, la 2 architettura e modularità dell'IA, la 3 sicurezza, privacy e conformità GDPR, la 4 accuratezza dell'IA e dashboard dei log, la 5 copertura di test e dataset MOCK; LG-n indica l'intera sezione;
		\item \textbf{Analisi del dominio}: requisiti ricavati dal gruppo dallo studio del dominio applicativo, non esplicitati nelle altre fonti.
	\end{itemize}
Le linee guida della proponente sono successive al verbale esterno: in caso di indicazioni in contrasto tra le fonti, il gruppo ha adottato quanto stabilito dalle linee guida.
